\section{The active best-seed branch as a phase--corner machine}

We keep the theorem object minimal. On the active best-seed branch we work with
\[
B=(s,u,v,\lambda,\mathrm{family}),\qquad \tau,\qquad \epsilon_4,
\]
and treat the auxiliary coordinates
\[
\rho=u_{\mathrm{source}}+1,
\qquad
\alpha=\rho-u,
\qquad
\delta=\rho-(s+u+v+\lambda),
\qquad
\kappa=q+s+v+\lambda
\]
as proof or constructive coordinates only.

\begin{theorem}[Phase--corner theorem]
Fix odd $m\ge 5$ and the best seed
\[
L=[2,2,1],\qquad R=[1,4,4].
\]
On the active nonterminal branch for the unresolved channel $R1\to H_{L1}$,
define
\[
\kappa=q+s+v+\lambda\pmod m.
\]
Then:
\begin{enumerate}
    \item $\kappa'\equiv \kappa+1\pmod m$.
    \item The current event $\epsilon_4$ is determined by $\kappa$, the carry
    sheet $c=1_{\{q=m-1\}}$, and the single corner state $s=2$ by the rules
    \begin{align*}
    \kappa=0 &\Rightarrow \mathrm{wrap},\\
    \kappa=1 &\Rightarrow \mathrm{carry\_jump},\\
    \kappa=2 &\Rightarrow
    \begin{cases}
    \mathrm{other}_{1000},& c=1,\\
    \mathrm{flat},& c=0,
    \end{cases}\\
    \kappa=3 &\Rightarrow
    \begin{cases}
    \mathrm{other}_{0010},& (c=0\ \&\ s=2)\ \text{or}\ (c=1\ \&\ s\ne 2),\\
    \mathrm{flat},& \text{otherwise},
    \end{cases}
    \end{align*}
    and for every $4\le \kappa\le m-1$ the event is flat.
    \item Among flat states, the only short reset occurs at the corner
    $(\kappa,s)=(2,2)$.
\end{enumerate}
Equivalently, the active branch is an odometer with one corner.
\end{theorem}

\begin{corollary}[Countdown carrier]
Define $\tau$ on the active branch by
\begin{itemize}
    \item $\tau=0$ on nonflat states;
    \item on flat states,
    \[
    \tau=
    \begin{cases}
    1,& (\kappa,s)=(2,2),\\
    m-\kappa,& \text{otherwise}.
    \end{cases}
    \]
\end{itemize}
Then $\tau'=\tau-1$ whenever $\tau>0$.
\end{corollary}

\begin{corollary}[Boundary reset formulas]
At $\tau=0$ the reset laws are:
\begin{itemize}
    \item $\mathrm{wrap}\mapsto 0$;
    \item $\mathrm{carry\_jump}\mapsto 0$ on the zero-reset fiber
    $s+v+\lambda\equiv 2\pmod m$, maps to $1$ when $s=1$ off that fiber, and to
    $m-2$ otherwise;
    \item $\mathrm{other}_{1000}\mapsto m-3$ if $s=1$ and $0$ otherwise;
    \item $\mathrm{other}_{0010}\mapsto m-4$.
\end{itemize}
\end{corollary}

\begin{corollary}[Finite-cover compatibility]
The phase--corner machine is compatible with the checked finite-cover normal
form
\[
B\leftarrow B+c\leftarrow B+c+d,
\]
where
\[
B=(s,u,v,\lambda,\mathrm{family}),
\qquad
c=1_{\{q=m-1\}},
\qquad
d=1_{\{\text{next carry }u\ge m-3\}}.
\]
\end{corollary}

\begin{proof}[Proof spine]
The clean proof spine is:
\[
033\to 062\to 059.
\]
First, one imports the exact trigger family $H_{L1}$ from the branch analysis.
Second, using the phase-$1$ source-residue invariant, one derives the
universal first-exit targets
\[
T_{\mathrm{reg}}=(m-1,m-2,1),
\qquad
T_{\mathrm{exc}}=(m-2,m-1,1).
\]
Third, before those targets are reached, the active branch cannot hit any of
the six modified current classes, so every pre-exit current state is
$B$-labelled. On those $B$-states the mixed witness rule applies exactly, and
in the current coordinate $\kappa=q+s+v+\lambda$ it yields the scheduler stated
in the theorem. Restricting the scheduler to flat states shows that the only
short reset occurs at $(\kappa,s)=(2,2)$, and the countdown and branchwise
reset formulas are immediate corollaries.
\end{proof}

\begin{remark}[Constructive refinement]
On the checked range through $m=19$, the stronger source-residue gauge with
\[
\rho=u_{\mathrm{source}}+1,
\qquad
\alpha=\rho-u,
\qquad
J=1_{\{\epsilon_4=\mathrm{carry\_jump}\}}
\]
gives exact current formulas
\[
q\equiv -\alpha+J\pmod m,
\qquad
c=1_{\{\alpha=1+J\}}.
\]
This is a constructive refinement of the theorem-side machine, not a change in
the theorem object.
\end{remark}
